\chapter{Estado del Arte}

\section{Secadores solares tipo invernadero en la literatura}



Los secadores solares constituyen una alternativa para el procesamiento de productos agrícolas y marinos. La revisión de \textcite{fudholi2010} distingue configuraciones directas, indirectas, mixtas e híbridas, y examina su desempeño técnico y económico. Esta diversidad exige relacionar las ventajas de cada diseño con el producto, las condiciones ambientales y los mecanismos de circulación utilizados.\\

Para el secado de café, la selección de una tecnología requiere examinar conjuntamente captación energética, transporte de vapor, uniformidad y costo de operación. Las posibilidades generales descritas en la literatura motivan el estudio de sistemas instrumentados, mientras que las prestaciones de una instalación particular deben establecerse mediante ensayos y condiciones de comparación documentados.\\

La certificación orgánica constituye una dimensión de la sostenibilidad de las fincas cafetaleras. \textcite{marquez2016} evaluaron indicadores de sustentabilidad ambiental antes y después de la certificación en La Convención, Cusco, Perú. Este antecedente sitúa la evaluación ambiental en un marco más amplio que debe complementarse con indicadores específicos del proceso poscosecha y del secado.\\

%Como es la situacion en Panama

En el caso panameño de Coffee Solutions, la variabilidad ambiental y la necesidad de caracterizar el microclima interior motivan la metodología de observación, modelado y diseño de control presentada en esta tesis. La pertinencia de automatizar el recinto se evalúa mediante variables y restricciones explícitas, conservando el alcance experimental de cada etapa.\\

En Moyobamba, Perú, \textcite{davila2022} desarrollaron un prototipo de secador solar tipo túnel para café Honey que incorporó elementos electrónicos, extracción de humedad y remoción automática del grano cada veinte minutos. El artículo reporta una temperatura máxima de $35\,^{\circ}\mathrm C$, humedad final entre $12$ y $14\%$ y una puntuación sensorial media de $82$ puntos en las muestras evaluadas. Estas cifras describen las condiciones y resultados de ese prototipo.\\

El antecedente muestra una aplicación concreta de instrumentación y actuación al secado solar. Su transferencia a otra instalación requiere caracterizar geometría, circulación, meteorología y producto, y evaluar por separado comportamiento del aire y calidad del café. Esta necesidad de correspondencia entre configuración y evidencia orienta la integración metodológica del presente trabajo.\\

La presente investigación está dedicada a establecer los fundamentos esenciales para el diseño y control de plantas de secado de café, con el objetivo primordial de lograr la estandarización del proceso de secado del café en Panamá. Este esfuerzo busca resolver la falta de protocolos definidos en el secado del café, que actualmente afecta la calidad sensorial del producto debido a la variabilidad en las condiciones climáticas y altitudinales. Al desarrollar y optimizar sistemas de secado avanzados, se pretende garantizar una calidad constante del café independientemente de la altitud, mejorando así la consistencia y el valor del producto final.\\

Además, la investigación se alinea con los objetivos de desarrollo sostenible al fomentar prácticas de secado más eficientes y sostenibles. Esto no solo contribuirá a la mejora de la calidad del café y la productividad, sino que también promoverá la sostenibilidad ambiental y económica en la producción de café. Al proporcionar una base sólida para el diseño y control de plantas, la investigación tiene el potencial de transformar la industria del café, adaptando el proceso de secado a las necesidades específicas de cada región y contribuyendo al desarrollo de prácticas más sostenibles en el sector.\\

Las referencias citadas en este capítulo se incluyen en la bibliografía general del documento.




\section{Técnicas modernas de simulación (CFD, híbridos)}

\section{Sistemas de control en secadores solares}

\section{Sensado y automatización en agricultura}

\section{Brechas detectadas y contribuciones del presente trabajo}
